Let \(U_W\) and \(U_A\) be independent fair binary coordinates and let \(U_0\) be independent of them, with state space \(\mathcal J_0=\{1,2,3,4\}\) and baseline masses \(u_j=1/4\).  Define the recursive structural functions
\[
 W_\star=U_W,
 \qquad A_\star=U_A,
\]
and define the \(Y_\star\)-response indexed by \(U_0=j\) by
\[
 f_{Y_\star}(a,j)=
 \begin{cases}
 0,&j=1,\\
 1,&j=2,\\
 a,&j=3,\\
 1-a,&j=4.
 \end{cases}
\]
Define the selection response by \(f_S(w,y,j)=1\) for \(j\in\{1,2\}\) and \(f_S(w,y,j)=0\) for \(j\in\{3,4\}\).  Each structural function uses only a subset of its displayed parents in \(\mathcal G_\star\), and the only shared exogenous coordinate, \(U_0\), is shared by \(Y_\star\) and \(S\).  Hence this is a recursive semi-Markovian SCM compatible with \(\mathcal G_\star\).  Its exogenous coordinates are independent, their declared state spaces equal their positive supports, and every atom has mass at least \(1/4\).

Selection occurs exactly for types \(1,2\), so \(P_M(S=1)=1/2\).  Conditional on selection, \(W_\star\) and \(A_\star\) remain independent fair coins, while \(Y_\star=0\) and \(Y_\star=1\) each have conditional probability \(1/2\).  Therefore
\[
 q_S(w,a,y)=\frac18\quad\text{for all }(w,a,y)\in\{0,1\}^3,
 \qquad q_0(w)=\frac12.
\]
Thus the input belongs to \(\mathcal C_{1/16}(\mathcal G_\star,T_0)\), and the baseline triple belongs to \(\mathfrak R_{1/16,1/4}^{\circ}\) for that input.

Order the active types as above and put \(\beta=(0,0,1,-1)^{\mathsf T}\).  Integrating over the two fair observed-variable coordinates gives
\[
 a_+=(1,1,0,0)^{\mathsf T},
 \qquad b_0=b_1=\tfrac12\mathbf 1.
\]
For a selected full cell with \(y=0\), \(a_v=(1/4,0,0,0)^{\mathsf T}\); for one with \(y=1\), \(a_v=(0,1/4,0,0)^{\mathsf T}\).  Since \(u^{\mathsf T}a_+=1/2\) and \(u^{\mathsf T}a_v=1/16\), every selected cross-product row is a scalar multiple of \((1,-1,0,0)^{\mathsf T}\).  Consequently \(\beta\) is orthogonal to \(\mathbf 1\), every \(b_t\), and every selected cross-product row, so
\[
 \beta\in\Gamma_{(M,\mathcal R,U_0)}^\perp.
\]
For the cell \((x,y)=(1,1)\), the intervention-response row is
\[
 h_{1,1}=(0,1,1,0)^{\mathsf T},
 \qquad \beta^{\mathsf T}h_{1,1}=1.
\]
Thus \(h_{1,1}\notin\Gamma_{(M,\mathcal R,U_0)}\).

Taking \(\tau=1/8=(1/4)/(2\lVert\beta\rVert_\infty)\), Theorem \(\mathrm{thm{:}projective\text{-}nullspace\text{-}certificate}\) applies.  More explicitly, the two \(U_0\)-mass vectors are
\[
 u^+=(1/4,1/4,3/8,1/8),
 \qquad u^-=(1/4,1/4,1/8,3/8).
\]
Their selected joint laws and \(W_\star\)-margins remain the displayed \(q_S,q_0\), while
\[
 \begin{aligned}
 P_{M_1}(Y_\star=1\mid\operatorname{do}(A_\star=1))
 &=u^+_2+u^+_3=\frac58,\\
 P_{M_2}(Y_\star=1\mid\operatorname{do}(A_\star=1))
 &=u^-_2+u^-_3=\frac38.
 \end{aligned}
\]
If \(\operatorname{EXT\text{-}ID}\) returned an arithmetic functional at this input, Theorem \(\mathrm{thm{:}extid\text{-}sound\text{-}termination}\) would force that single functional to equal both \(5/8\) and \(3/8\), a contradiction.  Termination leaves \(\operatorname{FAIL}\) as the only remaining output.  This proves the affirmative gate and its fully interior exact-fiber witness.